\section{Finite reproduction and the scope of the estimates}\label{dt:reproduction}

The exact finite identities and numerical budgets have a combined
reproduction command. From the repository root, run
\begin{verbatim}
python3 -I papers/circle-tree-volume/check.py \
    --output /tmp/circle-tree-volume-check
\end{verbatim}
The runner delegates to the two retained suites, in
\path{papers/circle-tree-asymptotics/check.py} and
\path{papers/directed-tree-asymptotics/check.py}. It records both
results together with the sources of this paper; the implementations
of the checks are not copied into a third suite.
The first script checks four symbolic edge identities, finite complex
products, and the summed-condition budgets. Its rational enclosures
use convergent series for \(\pi\), \(\sinh\) and \(\cosh\), and its
negative controls reject altered budgets. The second checks the
directed rational inequalities, endpoint pairings, and finite
recursions against full root-grid evaluations with first normals.
The recursion comparisons use formal independent vertex parameters;
they test the finite algebra and do not establish the existence of
additional knot diagrams. Its negative controls include an incorrect
outgoing condition at the final root and an odd stack of added leaves.
Both records retain the hashes of the sources they check.

These finite calculations do not prove the uniform complex estimates,
the Fourier-tail bounds, or the geometric identification. The analytic
arguments are written in Sections~\ref{cta:edge}--\ref{dt:sec-7};
the geometric input and its quantitative hypotheses are stated where
each theorem uses them. Background for the residue and Gaussian
calculations can be found in \cite[Sections~2.4(iii)--(iv)]{DLMF}.
The contours, endpoint terms and strict losses needed for the finite
invariant are specified in the proofs above.

The results concern fixed literal ordered circle trees. In particular,
they require the endpoint condition of Section~\ref{cta:intro}, and do
not apply solely from the abstract graph. The degree-and-twist
conditions do not address a tree that grows with \(N\). Neither proof
extends a short-twist subtree to an arbitrary multi-boundary tangle:
that would require additional conditional estimates for its boundary
variables. These are separate questions from the two volume theorems
proved here.
